% titel: Gekoppelte Pendel
% skript: Kap. Schwingungen

% Zwei gleiche Federpendel, durch Feder k gekoppelt; Entkopplung in
% Normalschwingungen (symmetrisch/antisymmetrisch), Energietransfer.

Zwei gleiche Oszillatoren, durch eine Koppelfeder verbunden: Die Bewegung ist
eine Überlagerung zweier Normalschwingungen, in denen beide Massen mit
derselben Frequenz schwingen.

\begin{werkzeugkasten}[Werkzeuge]
\textbf{Bewegungsgleichungen} (Kräfte der Koppelfeder beachten):
\begin{equation}
  m\,\ddot x_1=-D\,x_1+k(x_2-x_1),\qquad
  m\,\ddot x_2=-D\,x_2-k(x_2-x_1)\;.
\end{equation}

\textbf{Normalkoordinaten}:
\begin{equation}
  y_\pm\equiv\tfrac12(x_1\pm x_2)
  \;\Rightarrow\;
  \ddot y_++\omega_+^2 y_+=0,\quad
  \ddot y_--+\omega_-^2 y_-=0,
\end{equation}
mit $\omega_+=\sqrt{D/m}$ (gleichphasig, Koppelfeder ungedehnt) und
$\omega_-=\sqrt{(D+2k)/m}$ (gegenphasig).

\textbf{Lösung und Rücktransformation:}
\begin{equation}
  x_1=y_++y_-,\qquad x_2=y_+-y_-\;.
\end{equation}
Überlagerung zweier naher Frequenzen $\Rightarrow$ Schwebung: Die Energie
schwappt zwischen den Oszillatoren hin und her mit der Differenzfrequenz
$\omega_--\omega_+$.

\textbf{Einheiten:} $\omega_\pm$ in \unit{\per\second}, $D,k$ in
\unit{\newton\per\metre}.
\end{werkzeugkasten}

\begin{werkzeugkasten}[Vorgehen]
\begin{itemize}
  \item Gleichungen aufstellen: Federkräfte der \emph{eigenen} Feder $D\,x_i$
       plus Koppelfeder $k(x_2-x_1)$ mit Vorzeichen.
  \item Summe und Differenz der Gleichungen bilden $\Rightarrow$
       entkoppelte Gleichungen für $y_\pm$.
  \item Anfangsbedingungen auf $y_\pm(0)$, $\dot y_\pm(0)$ übertragen;
       Rücktransformieren liefert $x_1(t),x_2(t)$.
  \item Energietransfer: nach der Differenz der Normalfrequenzen fragen.
\end{itemize}
\end{werkzeugkasten}

\begin{achtung}[Typische Fallen]
\begin{itemize}
  \item Die Koppelfeder zieht an \emph{beiden} Massen: Kraft $+k(x_2-x_1)$
       an Masse 1, $+k(x_1-x_2)=-k(x_2-x_1)$ an Masse 2.
  \item $\omega_-\neq\omega_+$: Nur bei gleicher Masse und Feder sind beide
       Pendel identisch.
  \item Energietransfer-Frequenz ist $\omega_--\omega_+$, nicht $\omega_+$
       oder $\omega_-$ allein.
\end{itemize}
\end{achtung}
